\documentclass[10pt,letterpaper]{article}

\usepackage[left=0.5in,right=0.5in,top=0.5in,bottom=0.5in]{geometry}
\usepackage{enumitem}
\usepackage{hyperref}
\usepackage{titlesec}
\usepackage{mathptmx}

\pagestyle{empty}
\setlength{\parindent}{0pt}
\setlength{\parskip}{0pt}

\titleformat{\section}{\large\bfseries}{}{0em}{}[\titlerule]
\titlespacing*{\section}{0pt}{4pt}{2pt}

\setlist[itemize]{leftmargin=*,noitemsep,topsep=1pt,parsep=0pt,partopsep=0pt}

\begin{document}

\begin{center}
    {\LARGE \textbf{AMOGH BIJU}}\\[2pt]
    Los Angeles, California | 213-681-7507 | amogh@outlook.com | \href{https://www.linkedin.com/in/amogh-biju/}{linkedin.com/in/amogh-biju} | \href{https://github.com/gamefreakoneone}{github.com/gamefreakoneone}
\end{center}

\vspace{-10pt}

\section*{EDUCATION}
\textbf{University of Southern California} \hfill January 2024--December 2025\\
Master of Science in Computer Science

\vspace{-6pt}

\section*{EXPERIENCE}

\textbf{Research Assistant} \hfill February 2026--Present\\
\textit{USC Viterbi School of Engineering, Department of Computer Science, Los Angeles, CA}
\begin{itemize}
    \item Designed real-time fall detection system by fine-tuning Ultralytics YOLO11 on 111 annotated images, achieving 0.923 mAP@0.5, 94\% recall, and 0.90 F1 score for elderly safety monitoring
    \item Reduced false positive alerts by implementing a 3.5-second confirmation buffer with pose analysis, delivering email notifications with screenshots within 6--8 seconds of fall events
    \item Constructed a natural-language video query pipeline by combining SAM3 segmentation with temporal video analysis and integrating OpenAI and Gemini APIs, enabling queries over patient daily activities
\end{itemize}

\vspace{-2pt}

\textbf{Research Intern} \hfill June 2025--September 2025\\
\textit{USC Institute for Creative Technologies, Los Angeles, California}
\begin{itemize}
    \item Collaborated with team of 4 (2 PhD students, 1 postdoc, 1 peer) to build 3D reconstruction pipeline integrating RDD Structure-from-Motion, Doppelganger++ neural matching, and Gaussian Splatting
    \item Implemented feature extraction pipeline using RDD descriptors and LightGlue matcher, processing 500+ images per scene and reducing reconstruction time from 8 to 5 minutes (37\% improvement)
    \item Presented final pipeline demo to lab team; optimized Doppelganger++ duplicate detection through strategic image subset selection, improving point cloud quality
\end{itemize}

\vspace{-2pt}

\textbf{Research Intern, Department of Computer Science} \hfill June 2022--July 2023\\
\textit{IIT Hyderabad, Hyderabad, India}
\begin{itemize}
    \item Partnered with PhD researcher from Osaka University through weekly technical sessions, co-authoring IEEE Access survey paper analyzing 11 cellular network monitoring applications across 2G--5G architectures
    \item Captured physical layer cellular data using MobileInsight chipset-level monitoring and Dragonet's VPN-based approach, analyzing RRC/NAS signaling and traffic patterns across multiple Android devices
\end{itemize}

\vspace{-4pt}

\section*{PROJECTS}

\textbf{Defeat the Darkness -- AI-Powered Beat-Em-Up Game} \hfill \href{https://github.com/gamefreakoneone/all-your-base-is-ours}{\texttt{github.com/gamefreakoneone/all-your-base-is-ours}}
\begin{itemize}
    \item Constructed real-time multimodal AI game by integrating three Gemini 3 capabilities (image generation, structured output, live voice streaming), delivering unique procedurally-generated gameplay per session
    \item Architected parallel asset generation system by running 5 concurrent Gemini Pro image calls during player onboarding, eliminating perceived wait time despite multiple API round-trips
    \item Created self-correcting asset pipeline by implementing LLM-as-a-Judge validation with automatic retry and fallback, ensuring consistent sprite quality with a 70+ score threshold
\end{itemize}

\vspace{-2pt}

\textbf{Where We Left Off -- AI Reading Companion} \hfill \href{https://github.com/gamefreakoneone/Where-we-left-off-reader}{\texttt{github.com/gamefreakoneone/Where-we-left-off-reader}}
\begin{itemize}
    \item Developed spoiler-protected reading companion by implementing chapter-aware RAG with ChromaDB metadata filtering, enabling contextual Q\&A restricted to user's current reading progress
    \item Built agentic Q\&A system by designing LangGraph workflows with conditional semantic search and refinement loops, delivering contextually-grounded answers without future plot spoilers
    \item Crafted entity resolution pipeline by implementing alias canonicalization across chapters, enabling unified character tracking despite varying name references (e.g., Elizabeth/Lizzy/Ms.\ Bennet)
\end{itemize}

\vspace{-2pt}

\textbf{Ctrl+Alt+Del Hate -- Explainable Hate Speech Detection} \hfill \href{https://github.com/gamefreakoneone/Ctrl-Alt-Del-Hate}{\texttt{github.com/gamefreakoneone/Ctrl-Alt-Del-Hate}}
\begin{itemize}
    \item Trained multi-task hate speech detection model by LoRA-finetuning LLaMA-3.2-1B on 39K annotated comments, achieving 83.7\% accuracy and 0.693 F1 on multi-class classification
    \item Formulated three-task classification architecture with shared encoder and task-specific heads, delivering 0.59 Micro F1 on 45 demographic target groups and 0.53 MAE on 10 ordinal toxicity facets
    \item Attained 0.80 Spearman correlation on hate speech severity scoring, enabling fine-grained content moderation beyond binary classification
\end{itemize}

\vspace{-4pt}

\section*{TECHNICAL SKILLS}
\textbf{Languages:} Python, JavaScript, TypeScript, SQL, C++, HTML/CSS, Bash \\[2pt]
\textbf{AI/ML:} PyTorch, YOLO, LLaMA, LoRA, LangChain, LangGraph, Hugging Face Transformers, OpenAI API, Gemini API, SAM3, RAG \\[2pt]
\textbf{Web \& Backend:} React, Next.js, Flask, FastAPI, Node.js, REST APIs, WebSockets, Chrome Extensions, Vercel \\[2pt]
\textbf{Data \& Infrastructure:} MongoDB, PostgreSQL, ChromaDB, Redis, Docker, AWS (S3, EC2), Git, Linux \\[2pt]
\textbf{Computer Vision \& 3D:} OpenCV, COLMAP, Gaussian Splatting, Structure-from-Motion, LightGlue, RDD Descriptors

\end{document}
